% Fichier engendré par outils/generer-tex.py à partir de 06-integration.
% Les démonstrations en français sont transcrites du fichier Lean (skill transcrire-preuve-lean).
\documentclass[11pt,a4paper]{article}

% Compile aussi bien avec pdflatex qu'avec xelatex ou tectonic.
\usepackage{iftex}
\ifPDFTeX
  \usepackage[T1]{fontenc}
  \usepackage[utf8]{inputenc}
\else
  \usepackage{fontspec}
\fi
\usepackage[french]{babel}
\usepackage{amsmath,amssymb,amsthm}
\usepackage[margin=2.5cm]{geometry}
\usepackage[hidelinks]{hyperref}

\theoremstyle{definition}
\newtheorem{definition}{Définition}
\theoremstyle{plain}
\newtheorem{theoreme}{Théorème}
\newtheorem{lemme}[theoreme]{Lemme}

% Renvoi à la déclaration Lean, une ligne après la démonstration, aligné à droite.
\newcommand{\source}[2]{\par\smallskip\noindent\hfill{\footnotesize\href{#1}{\texttt{#2}}}\par\smallskip}

\title{Intégration}
\date{}

\begin{document}
\maketitle
% Les identifiants Lean en machine à écrire ne se coupent pas : on tolère des
% espacements plus lâches plutôt que des lignes qui débordent.
\sloppy

\section{Integration}

\noindent
Lycée \textemdash{} section \og{} Intégration \fg{}, niveau terminale S.

\medskip\noindent
Le programme \emph{définit} l'intégrale d'une fonction continue positive comme l'aire
sous la courbe, et admet que cette aire existe. Mathlib fait l'inverse : l'intégrale y est
construite par la théorie de la mesure, et l'aire est la mesure d'une partie du plan. Les
deux points de vue ne se rejoignent pas par convention mais par un théorème, qui ouvre ce
chapitre : la mesure de la partie du plan comprise entre l'axe et la courbe \emph{vaut}
l'intégrale. C'est l'énoncé le plus intéressant du fichier, parce que c'est celui que le
lycée ne peut qu'admettre.

\medskip\noindent
Notation : $\int_a^b f$ désigne l'intégrale sur l'intervalle orienté de $a$ à $b$, qui
change de signe quand on échange les bornes. Toutes les fonctions considérées sont
continues, ce qui rend l'intégrabilité automatique.

\subsection{L'intégrale est l'aire sous la courbe}

\begin{theoreme}
Soit $f$ continue et positive et $a \le b$. La partie du plan comprise entre l'axe des
abscisses et la courbe de $f$, au-dessus de $]a\,;b]$, a pour aire
\[ \lambda_2\Bigl(\bigl\{(x, y) \ :\ a < x \le b,\ 0 < y < f(x)\bigr\}\Bigr)
   = \int_a^b f(x)\,dx , \]
où $\lambda_2$ est la mesure de Lebesgue du plan.
\end{theoreme}

\begin{proof}
Mathlib appelle \emph{région entre deux courbes} l'ensemble
\[ \{(x, y) \ :\ x \in s,\ g(x) < y < h(x)\} , \]
et démontre que sa mesure vaut $\int_s (h - g)$ dès que $g$ et $h$ sont intégrables sur
$s$, que $s$ est mesurable et que $g \le h$ sur $s$. C'est le théorème de Fubini appliqué
à cette partie du plan : on mesure d'abord chaque tranche verticale, de longueur
$h(x) - g(x)$, puis on intègre ces longueurs.

On l'applique ici avec $g = 0$, $h = f$ et $s = \,]a\,;b]$ : la fonction nulle est
intégrable, $f$ l'est parce qu'elle est continue sur un intervalle borné, l'intervalle
$]a\,;b]$ est mesurable, et $0 \le f$ par hypothèse. Il reste
\[ \lambda_2 = \int_{]a;b]} (f - 0) = \int_{]a;b]} f , \]
et comme $a \le b$, l'intégrale sur l'intervalle orienté coïncide avec l'intégrale sur
$]a\,;b]$.

\medskip\noindent
Trois remarques sur l'écart avec l'énoncé du lycée. La partie mesurée est ouverte en $y$
et semi-ouverte en $x$ : ces bords sont de mesure nulle et ne changent pas l'aire, mais
c'est la forme sous laquelle la bibliothèque énonce le résultat. La mesure est à valeurs
dans $[0, +\infty]$, d'où la conversion explicite du réel $\int_a^b f$ en un élément de cet
ensemble. Enfin, ce théorème est l'exact renversement du cours : au lycée l'aire définit
l'intégrale, ici l'intégrale calcule l'aire.
\end{proof}
\source{https://github.com/Commutator-IO/learn-lean/blob/main/cours/02-lycee/06-integration/Integration.lean\#L21}{Integration.lean\#L21}

\subsection{Primitives}

\begin{theoreme}
Toute fonction continue admet des primitives, et
\[ F(x) = \int_a^x f(t)\,dt \]
est celle qui s'annule en $a$ : $F' = f$ et $F(a) = 0$.
\end{theoreme}

\begin{proof}
Le fait que $F$ soit dérivable de dérivée $f$ est le théorème fondamental de l'analyse,
démontré dans la bibliothèque pour toute fonction continue : le taux d'accroissement de
$F$ entre $x$ et $x + h$ vaut la valeur moyenne de $f$ sur $[x\,;x+h]$, qui tend vers
$f(x)$ par continuité. La bibliothèque en donne même la version forte (dérivée stricte),
dont on ne retient que la dérivabilité ordinaire.

L'annulation en $a$ est immédiate : une intégrale dont les deux bornes sont égales est
nulle.

\medskip\noindent
Le programme admet l'existence de primitives ; ici elle est \emph{construite}, et la
formule qui la donne est celle de l'énoncé suivant du programme. Les deux points ne font
donc qu'un.
\end{proof}
\source{https://github.com/Commutator-IO/learn-lean/blob/main/cours/02-lycee/06-integration/Integration.lean\#L34}{Integration.lean\#L34}

\begin{theoreme}
Deux primitives d'une même fonction diffèrent d'une constante : si $F' = G' = f$ sur
$\mathbb{R}$, il existe $c$ tel que $F(x) - G(x) = c$ pour tout $x$.
\end{theoreme}

\begin{proof}
Il suffit de prendre $c = F(0) - G(0)$.

La fonction $t \mapsto F(t) - G(t)$ est dérivable, de dérivée
\[ F'(y) - G'(y) = f(y) - f(y) = 0 \]
en tout point $y$. Une fonction dérivable sur $\mathbb{R}$ de dérivée identiquement nulle
est constante \textemdash{} conséquence du théorème des accroissements finis, démontrée
dans la bibliothèque \textemdash{} donc sa valeur en $x$ est celle en $0$, ce qui est
exactement l'égalité annoncée.
\end{proof}
\source{https://github.com/Commutator-IO/learn-lean/blob/main/cours/02-lycee/06-integration/Integration.lean\#L39}{Integration.lean\#L39}

\subsection{Théorème fondamental de l'analyse}

\begin{theoreme}
Si $f$ est continue et si $F$ est une primitive de $f$ sur $\mathbb{R}$, alors
\[ \int_a^b f(x)\,dx = F(b) - F(a) . \]
\end{theoreme}

\begin{proof}
C'est le théorème de la bibliothèque : si $F$ est dérivable en tout point du segment
d'extrémités $a$ et $b$, de dérivée $f$, et si $f$ y est intégrable, alors l'intégrale de
$f$ vaut l'accroissement de $F$. Les deux hypothèses sont ici satisfaites partout :
$F' = f$ sur $\mathbb{R}$ tout entier, et $f$ continue est intégrable sur tout segment.

\medskip\noindent
L'énoncé ne suppose pas $a \le b$ : l'intégrale sur l'intervalle orienté change de signe
avec l'ordre des bornes, exactement comme $F(b) - F(a)$.
\end{proof}
\source{https://github.com/Commutator-IO/learn-lean/blob/main/cours/02-lycee/06-integration/Integration.lean\#L55}{Integration.lean\#L55}

\subsection{Linéarité}

\begin{theoreme}
\[ \int_a^b \bigl(f(x) + g(x)\bigr)\,dx = \int_a^b f(x)\,dx + \int_a^b g(x)\,dx ,
   \qquad \int_a^b k\,f(x)\,dx = k \int_a^b f(x)\,dx . \]
\end{theoreme}

\begin{proof}
Les deux règles sont celles de la bibliothèque. L'additivité réclame que $f$ et $g$ soient
intégrables sur l'intervalle, ce que leur continuité assure ; l'homogénéité ne demande
rien, la constante sortant de l'intégrale même lorsque celle-ci n'est pas définie.
\end{proof}
\source{https://github.com/Commutator-IO/learn-lean/blob/main/cours/02-lycee/06-integration/Integration.lean\#L62}{Integration.lean\#L62}

\subsection{Relation de Chasles}

\begin{theoreme}
\[ \int_a^b f(x)\,dx + \int_b^c f(x)\,dx = \int_a^c f(x)\,dx . \]
\end{theoreme}

\begin{proof}
Résultat de la bibliothèque, sous l'hypothèse que $f$ est intégrable sur les deux
premiers intervalles \textemdash{} vrai ici par continuité.

\medskip\noindent
Aucune condition d'ordre sur $a$, $b$, $c$ : c'est justement l'intérêt de l'intégrale
orientée, qui rend la relation vraie quelle que soit la position relative des trois
bornes, y compris lorsque $b$ n'est pas entre $a$ et $c$.
\end{proof}
\source{https://github.com/Commutator-IO/learn-lean/blob/main/cours/02-lycee/06-integration/Integration.lean\#L72}{Integration.lean\#L72}

\subsection{Positivité et croissance}

\begin{theoreme}
Pour $a \le b$ :
\[ f \ge 0 \ \text{sur } [a\,;b] \implies \int_a^b f \ge 0 , \qquad
   f \le g \ \text{sur } [a\,;b] \implies \int_a^b f \le \int_a^b g . \]
\end{theoreme}

\begin{proof}
Deux résultats de la bibliothèque, cités. La positivité vient de ce que l'intégrale d'une
fonction positive est une limite de sommes de termes positifs ; la croissance s'en déduit
en appliquant la positivité à $g - f$.

\medskip\noindent
L'hypothèse $a \le b$ est indispensable dans les deux cas : sur l'intervalle orienté à
l'envers, l'intégrale d'une fonction positive est négative.
\end{proof}
\source{https://github.com/Commutator-IO/learn-lean/blob/main/cours/02-lycee/06-integration/Integration.lean\#L81}{Integration.lean\#L81}

\subsection{Inégalité de la moyenne}

\begin{theoreme}
Si $m \le f \le M$ sur $[a\,;b]$ avec $a \le b$, alors
\[ m\,(b - a) \le \int_a^b f(x)\,dx \le M\,(b - a) , \]
et si de plus $a < b$, la valeur moyenne de $f$ vérifie
\[ m \le \frac{1}{b - a}\int_a^b f(x)\,dx \le M . \]
\end{theoreme}

\begin{proof}
On compare $f$ aux deux fonctions constantes. La croissance de l'intégrale donne
\[ \int_a^b m\,dx \le \int_a^b f(x)\,dx \le \int_a^b M\,dx , \]
et l'intégrale d'une constante sur $[a\,;b]$ vaut cette constante multipliée par la
longueur :
\[ \int_a^b m\,dx = (b - a)\,m , \qquad \int_a^b M\,dx = (b - a)\,M . \]
C'est l'encadrement annoncé, aux facteurs près écrits dans l'autre ordre.

Pour la valeur moyenne, on suppose $a < b$, donc $b - a > 0$, et l'on divise l'encadrement
par cette quantité strictement positive, ce qui en conserve le sens.

\medskip\noindent
Géométriquement : l'aire sous la courbe est comprise entre celles des deux rectangles de
base $[a\,;b]$ et de hauteurs $m$ et $M$, et la valeur moyenne est la hauteur du rectangle
de même aire.
\end{proof}
\source{https://github.com/Commutator-IO/learn-lean/blob/main/cours/02-lycee/06-integration/Integration.lean\#L94}{Integration.lean\#L94}

\subsection{Intégration par parties}

\begin{theoreme}
Si $u$ et $v$ sont dérivables sur $\mathbb{R}$, de dérivées $u'$ et $v'$ continues, alors
\[ \int_a^b u(x)\,v'(x)\,dx
   = u(b)v(b) - u(a)v(a) - \int_a^b u'(x)\,v(x)\,dx . \]
\end{theoreme}

\begin{proof}
Résultat de la bibliothèque, cité. Il découle du théorème fondamental appliqué au produit
$uv$ : la dérivée de $uv$ est $u'v + uv'$, donc
\[ \int_a^b \bigl(u'v + uv'\bigr) = u(b)v(b) - u(a)v(a) , \]
et il ne reste qu'à séparer les deux termes par linéarité. Les hypothèses de continuité
des dérivées servent à cette séparation, en assurant que chacune des deux intégrales
existe séparément.
\end{proof}
\source{https://github.com/Commutator-IO/learn-lean/blob/main/cours/02-lycee/06-integration/Integration.lean\#L119}{Integration.lean\#L119}

\subsection{Aire entre deux courbes}

\begin{theoreme}
Si $f \le g$ sont continues et $a \le b$, la partie du plan comprise entre les deux
courbes au-dessus de $]a\,;b]$ a pour aire
\[ \lambda_2\Bigl(\bigl\{(x, y) \ :\ a < x \le b,\ f(x) < y < g(x)\bigr\}\Bigr)
   = \int_a^b \bigl(g(x) - f(x)\bigr)\,dx . \]
\end{theoreme}

\begin{proof}
C'est le théorème qui ouvre le chapitre, appliqué cette fois avec $f$ et $g$ quelconques
au lieu de la fonction nulle et de $f$ : les deux fonctions sont intégrables sur
$]a\,;b]$ par continuité, l'intervalle est mesurable, et $f \le g$ par hypothèse. La
mesure vaut donc $\int_{]a;b]} (g - f)$, que l'on récrit sur l'intervalle orienté puisque
$a \le b$.

\medskip\noindent
\textbf{Ce qui n'est pas démontré.} Le volume d'un solide de révolution,
\[ V = \pi \int_a^b f(x)^{2}\,dx , \]
n'est pas formalisé. Il faudrait mesurer dans $\mathbb{R}^{3}$ la partie
\[ \bigl\{(x, y, z) \ :\ a \le x \le b,\ y^{2} + z^{2} \le f(x)^{2}\bigr\} , \]
donc disposer du théorème de Fubini en dimension trois et de l'aire du disque \textemdash{}
un appareillage sans commune mesure avec l'approche du programme, qui empile des cylindres
et passe à la limite sans le justifier. L'énoncé correspondant reste marqué
\og{} en cours \fg{} dans l'index du chapitre.
\end{proof}
\source{https://github.com/Commutator-IO/learn-lean/blob/main/cours/02-lycee/06-integration/Integration.lean\#L131}{Integration.lean\#L131}

\vfill
\noindent\rule{0.35\linewidth}{0.4pt}\par
\noindent{\footnotesize Leonardo de Moura et Sebastian Ullrich,
\emph{The Lean~4 Theorem Prover and Programming Language},
\emph{Automated Deduction — CADE~28}, Springer, 2021, p.~625--635,
\href{https://doi.org/10.1007/978-3-030-79876-5_37}{doi:10.1007/978-3-030-79876-5\_37}.\par}

\end{document}
